\documentclass[11pt]{article}
\usepackage[a4paper,margin=25mm]{geometry}
\usepackage{amsmath,amssymb,amsthm}
\usepackage[hidelinks]{hyperref}
\newtheorem{theorem}{Theorem}
\newtheorem{lemma}{Lemma}
\newcommand{\Kzero}{K_0}
\newcommand{\Fp}{\mathbb F}
\title{Different Fields Have the Same Scaled Ordered $K_0$:\\
A Disproof of Conjecture 00000006291}
\author{GitHub: gaochengzhecpu}
\date{}
\begin{document}
\maketitle
\begin{abstract}
The assertion that the zeroth $K$-group completely classifies regular rings
fails already for the fields $\Fp_2$ and $\Fp_3$. Both are von Neumann regular,
and their scaled ordered $K_0$ groups are isomorphic to
$(\mathbb Z,\mathbb Z_{\geq0},1)$, although the rings are not isomorphic.
The Lean proof constructs the Grothendieck group from actual finite-dimensional
vector spaces and linear equivalences, verifies its universal property, and
proves the counterexample with the positive cone and the unit class retained.
\end{abstract}

\paragraph{Reading of the source.}
Both language versions assert that the zeroth group is a complete invariant
of regular rings and describe a classification by ordered groups.
We use the standard unital von Neumann meaning of regular:
\[
 \forall a\in R\ \exists b\in R,\qquad aba=a.
\]
The zeroth group is the algebraic Grothendieck group $K_0(R)$ of finitely
generated projective right $R$-modules under direct sum. Its positive cone
consists of the classes of actual modules, and its distinguished order unit
is $[R]$. The rings here are commutative, so right versus left is immaterial.
Completeness for ring isomorphism would imply, in particular,
\begin{equation}\label{eq:claim}
 (K_0(F),K_0(F)^+,[F])\cong(K_0(E),K_0(E)^+,[E])
 \quad\Longrightarrow\quad F\cong E
\end{equation}
for every pair of fields $F,E$, since fields are regular.
We disprove this necessary restriction. Retaining the positive cone and
unit makes the counterexample stronger than one involving only abstract
abelian groups. No common ground field or narrower class of regular rings
is stipulated in the source; a classification theorem with such extra
restrictions is outside the claim refuted here.

\begin{theorem}\label{thm:main}
The rings $\Fp_2$ and $\Fp_3$ are von Neumann regular and nonisomorphic, but
\[
 (K_0(\Fp_2),K_0(\Fp_2)^+,[\Fp_2])
 \cong(\mathbb Z,\mathbb Z_{\geq0},1)
 \cong(K_0(\Fp_3),K_0(\Fp_3)^+,[\Fp_3]).
\]
Consequently \eqref{eq:claim}, and hence the stated complete-invariant
assertion for regular rings, is false.
\end{theorem}

\begin{lemma}\label{lem:field}
For every field $F$, dimension induces a scaled ordered group isomorphism
\[
 (K_0(F),K_0(F)^+,[F])\cong(\mathbb Z,\mathbb Z_{\geq0},1).
\]
\end{lemma}
\begin{proof}
Every finitely generated $F$-module is a finite-dimensional vector space
and has a finite basis. A free module is projective: to lift a linear map
through a surjection, choose preimages of the images of basis vectors and
extend linearly. Thus the finitely generated projective modules over $F$
are precisely its finite-dimensional vector spaces. Their isomorphism
classes form a commutative monoid under finite direct sum, which can be
realized as the Cartesian product with componentwise vector operations.
Dimension identifies this monoid with $(\mathbb N,+)$, because
\[
 P\cong F^{\dim_F P},\qquad
 P\cong Q\ \Longleftrightarrow\ \dim_F P=\dim_F Q,
 \qquad \dim_F(P\oplus Q)=\dim_F P+\dim_F Q.
\]
The following explicit group completion also specifies the construction
used in Lean.

Take pairs $(P,Q)$ of actual finite-dimensional vector spaces. Declare
\begin{equation}\label{eq:rel}
 (P,Q)\sim(P',Q')
 \quad\Longleftrightarrow\quad
 P\oplus Q'\cong P'\oplus Q.
\end{equation}
This is equivalent to equality of the integer differences
$\dim P-\dim Q$ and $\dim P'-\dim Q'$ and is therefore an equivalence
relation. It also agrees with the customary stabilized relation allowing
an auxiliary summand $T$ on both sides of the isomorphism in
\eqref{eq:rel}: equality of the resulting dimensions cancels $\dim T$,
and the converse follows by taking $T=0$.

On equivalence classes put
\[
 [(P,Q)]+[(P',Q')]=[(P\oplus P',Q\oplus Q')],
 \qquad -[(P,Q)]=[(Q,P)].
\]
These operations are well-defined by the dimension criterion. The map
\[
 d([(P,Q)])=\dim_F P-\dim_F Q
\]
is well-defined and injective by \eqref{eq:rel}, and is surjective because
an integer $z$ is represented by $(F^{\max(z,0)},F^{\max(-z,0)})$.
It is an additive isomorphism to $\mathbb Z$.

For completeness, this quotient has the defining universal property of
the Grothendieck group. Let $A$ be any abelian group and let $f$ assign an
element of $A$ to each finite-dimensional vector space, invariant under
linear isomorphism, additive under direct sum, and zero on the zero space.
Then
\[
 [(P,Q)]\longmapsto f(P)-f(Q)
\]
is well-defined: an isomorphism in \eqref{eq:rel} gives
$f(P)+f(Q')=f(P')+f(Q)$. This map is additive and sends $[(P,0)]$ to
$f(P)$. It is unique, because every class equals
$[(P,0)]-[(Q,0)]$. Hence the quotient is the Grothendieck group of the
specified monoid, namely $K_0(F)$. Equivalently, the exact-sequence
presentation gives the same group: short exact sequences of projective
modules split, and their relations are exactly direct-sum relations.

The positive cone is the set of classes $[(P,0)]$. Their images under $d$
are precisely the nonnegative integers, since $\dim_F P\geq0$ and every
$n\geq0$ occurs for $P=F^n$. Thus the order is exactly
\[
 x\leq y\quad\Longleftrightarrow\quad y-x\in K_0(F)^+
 \quad\Longleftrightarrow\quad d(x)\leq d(y).
\]
Finally $[F]$ has dimension one. This proves the scaled ordered assertion.
\end{proof}

\begin{proof}[Proof of Theorem~\ref{thm:main}]
In any field, take $b=a^{-1}$ when $a\neq0$ and $b=0$ when $a=0$;
then $aba=a$. Thus both fields are von Neumann regular.
Apply Lemma~\ref{lem:field} to the two fields and compose one dimension
isomorphism with the inverse of the other. The composition preserves
addition, order, and the distinguished unit. A ring isomorphism between
$\Fp_2$ and $\Fp_3$ would be a bijection between their underlying sets,
which have different cardinalities $2$ and $3$. No such isomorphism exists.
\end{proof}

\section*{Formal verification and scope}
The self-contained file \texttt{lean/Main.lean} uses Lean 4.19.0 and
Mathlib at commit
\begin{center}\texttt{c44e0c8ee63ca166450922a373c7409c5d26b00b}.\end{center}
Its bundled objects carry actual types, additive groups, scalar actions,
and Mathlib's finite-generation instances. The following bridges are
proved in Lean:
\begin{itemize}
\item Every bundled object is finitely generated and projective, every
finitely generated module can be bundled, and every object is linearly
isomorphic to the standard finite-dimensional space of its dimension.
\item The defining relation is a genuine linear equivalence between
cross direct sums; allowing an auxiliary stabilizing module gives the
same relation.
\item The rank map on the quotient is bijective. The abelian-group
structure is transported along it, then proved to be direct sum on pair
representatives, with coordinate exchange as additive inverse.
\item The group completion has its universal property for every abelian
target group. Its partial order is compatible with addition and is
proved to be exactly the order coming from actual module classes.
\item The unit class equals the class of the field as a module over
itself. For every pair of fields an ordered additive group isomorphism
preserving this unit is constructed.
\item The two \texttt{ZMod} fields are regular, no ring isomorphism exists
between them, and the universally quantified field-completeness
assertion is negated.
\end{itemize}
The final theorem is
\begin{center}\texttt{Conjecture6291.conjecture\_6291\_false}.\end{center}
Its type is the negation of \texttt{FieldCompleteness}, which quantifies
over every pair of fields and assumes regularity and an isomorphism of
their scaled ordered groups before demanding a ring isomorphism.
This is an explicitly stated necessary restriction of the source's
complete-invariant assertion, not a formal construction of $K_0$ for
arbitrary rings. The finite-projective/direct-sum definition is
formalized with its universal property; the equivalence with the
short-exact-sequence presentation is explained on paper. The remaining
unspecified claims concerning continuous geometry are unnecessary once
completeness has failed and are not asserted here. The axiom audits use
only \texttt{propext}, \texttt{Classical.choice}, and \texttt{Quot.sound}.

\section*{Source and reproduction}
The exact bilingual statement is included as \texttt{SOURCE.md}, from
upstream commit
\begin{center}\texttt{b355e42906e6eaf52bc7e9d005e0ef02d8aa7400}.\end{center}
Its source hash and retrieval record are in
\begin{center}\texttt{verification/SOURCE\_PROVENANCE.md}.\end{center}
The computation $K_0(F)\cong\mathbb Z$ by dimension is also recorded in
the Stacks Project, Section 10.55, Example 10.55.2:
\href{https://stacks.math.columbia.edu/tag/00JC}{\texttt{stacks.math.columbia.edu/tag/00JC}}.
The proof above supplies the group, cone, and unit calculation directly.

With the pinned dependencies available, run in the \texttt{lean} directory:
\begin{center}\texttt{lake env lean -DwarningAsError=true Main.lean}\end{center}
and \texttt{lake build} for a fresh project build. Compile
\texttt{main.tex} with Tectonic or a compatible LaTeX installation.
This submission was developed by a delegated AI agent with a self-review;
the coordinating agent reviews it separately. No external independent
review is claimed.
\end{document}
